% One worked example, one guided example and four independent practice
% problems cover the six core families once each, in source order.
% Matching guided and independent review solutions are optional.
\makeatletter
\def\input@path{{../../../preamble/}}
\makeatother
\input{preamble.tex}
\input{graphs/highlights.tex}
\GTenHighlightsfalse % Disable instructional highlights throughout this activity.
\tikzset{gten marker/.append style={inner xsep=.6mm,inner ysep=.3mm}}
% Local C2 roles follow C1: quantities orange; probabilities in related blues.
\colorlet{cTwoMassMarker}{GTenHLProbability}
\definecolor{cTwoCDFMarker}{HTML}{43A5C9}
\newcommand{\cTwoValue}[1]{\HighlightStroke{GTenHLValue}{#1}}
\newcommand{\cTwoOutcome}[1]{\HighlightStroke{GTenHLOutcome}{#1}}
\newcommand{\cTwoMassHL}[1]{\HighlightStroke{cTwoMassMarker}{#1}}
\newcommand{\cTwoCDFHL}[1]{\HighlightStroke{cTwoCDFMarker}{#1}}
\colorlet{cTwoMassColor}{blue!75!black}
\colorlet{cTwoCDFColor}{cyan!65!blue!80!black}
\colorlet{cTwoGuideColor}{cTwoMassColor!35}
\ifGTenHighlights\else
  \colorlet{cTwoMassColor}{black}
  \colorlet{cTwoCDFColor}{black}
  \colorlet{cTwoGuideColor}{gray}
\fi
\input{graphs/G10_T2_bar_graphs.tex}
\usepackage{needspace}
\usepackage{varwidth}
\geometry{letterpaper,margin=0.25in}
\pagestyle{empty}
\linespread{1.04}
\setlength{\parskip}{4pt}
\setlength{\tabcolsep}{6pt}
\setlength{\extrarowheight}{2pt}
\setlength{\emergencystretch}{2em}
\hypersetup{pdftitle={Grade 10 C2 Classroom Activity: PMF and CDF Tables and Graphs}}
\newif\ifprintsolutions
\printsolutionsfalse % Set \printsolutionstrue to append the complete review solutions.
\newcolumntype{Y}{>{\centering\arraybackslash}X}
\newcommand{\stageheading}[1]{\par\Needspace{6\baselineskip}\addvspace{.65\baselineskip}{\centering\large\bfseries #1\par}\nopagebreak[4]}
\newcommand{\problemtitle}[3]{\par\Needspace{4\baselineskip}\addvspace{.4\baselineskip}\noindent{\setlength{\fboxsep}{2pt}\colorbox{black!10}{\begin{varwidth}[t]{.75\linewidth}\bfseries\strut #1. #2\end{varwidth}}}\hfill\textbf{Type #3}\par\nopagebreak[4]}
\newcommand{\task}[2]{\par\noindent\textbf{Task #1 --- \ifnum#1=1 PMF table and graph\else CDF table and graph\fi.} #2\par}
\newcommand{\responseline}{\par\noindent\rule{0pt}{1.55em}\hrulefill\par}
\newcommand{\worklines}[1]{\begingroup\setlength{\parskip}{0pt}\foreach \i in {1,...,#1}{\responseline}\endgroup}
\newcommand{\blankgraph}[1]{\begin{tikzpicture}\begin{axis}[width=\linewidth,height=4.4cm,axis lines=left,xmin=0,xmax=10,ymin=0,ymax=1.05,xtick=\empty,ytick={0,.25,.5,.75,1},grid=major,grid style={gray!20},xlabel={$x$ (USD)},ylabel={#1},tick label style={font=\small},label style={font=\small}]\end{axis}\end{tikzpicture}}
\newcommand{\blanktable}[1]{\begin{tabularx}{\linewidth}{|Y|Y|}\hline #1\\\hline\rule{0pt}{.6cm}&\\\hline\rule{0pt}{.6cm}&\\\hline\rule{0pt}{.6cm}&\\\hline\rule{0pt}{.6cm}&\\\hline\rule{0pt}{.6cm}&\\\hline\rule{0pt}{.6cm}&\\\hline\end{tabularx}}
\pgfplotsset{classCDF/.style={width=\dimexpr\linewidth-6pt\relax,height=4.15cm,axis lines=left,ymin=-.06,ymax=1.1,ytick={0,.25,.5,.75,1},ylabel={$\cTwoCDFHL{F_X(x)}$},xlabel={$\cTwoValue{x}$ (\cTwoValue{USD})},grid=major,grid style={gray!15},tick label style={font=\small},label style={font=\small},clip=false}}
\begin{document}
{\centering{\Large\bfseries C2 -- Constructing PMF and CDF Tables and Graphs\par}
Global Economics \textbar{} Grade 10 \textbar{} Term 2 \textbar{} Lesson 1\par}
\textbf{C2:} Interprets PMFs and CDFs by verifying their defining properties and using them to calculate point, interval, cumulative, and complementary probabilities in financial contexts.

\stageheading{Equations and basic procedure}
\noindent\begin{tabularx}{\linewidth}{@{}l l X@{}}\toprule
Symbol or notation & Concept & Brief meaning\\\midrule
$\cTwoValue{X} \>\> | \>\> \cTwoValue{x} \>\> | \>\> \cTwoValue{S_X}$ & \cTwoValue{Quantity} & Variable | Value | Support of the random variable $X$\\
$\cTwoOutcome{\omega}$ & \cTwoOutcome{Outcome} & One complete result.\\
$\cTwoMassHL{p_X(x) = P(X=x)}$ & \cTwoMassHL{PMF} & The \textbf{Probability Mass Function (PMF)} describes the probability that a discrete random variable $X$ takes a specific value $x$. \\
$\cTwoCDFHL{F_X(x)= P(X \leq x)}$ & \cTwoCDFHL{CDF} & The \textbf{Cumulative Distribution Function (CDF)} describes the probability that a random variable $X$ takes a value less than or equal to a specified value $x$.\\\bottomrule
\end{tabularx}

\task{1}{Define the random variable clearly stating what it represents and its possible values or \textbf{Support}. Construct a PMF table showing the probability of each possible value, and draw a labeled graph representing the probability distribution.}
\task{2}{Construct a CDF table showing the cumulative probability for each interval of values, including the intervals below the minimum and at or above the maximum possible values. Draw a labeled step graph representing the cumulative distribution, using filled and open points to indicate included and excluded endpoints.}

\stageheading{Worked example} \vspace{-5mm}
\problemtitle{1}{Museum kiosk payments}{A}
\begingroup\small
\textbf{Payment fees.} A museum kiosk receives one payment. The only methods are cash, debit card, credit card and mobile wallet, with probabilities .15, .25, .35, .25, respectively. Cash has no fee; either card costs 2 USD; a mobile wallet costs 4 USD. There are no other charges. Record the processing fee.

\task{1}{Let $\cTwoValue{X}$ be the \cTwoValue{processing fee} for one payment, in \cTwoValue{USD}. The support is $S_X=\{\cTwoValue{0},\cTwoValue{2},\cTwoValue{4}\}$.}
% Separate top-aligned blocks keep tables and graph at their native text size.
\par\noindent\begin{minipage}[t]{.28\linewidth}\vspace{0pt}
\small\setlength{\tabcolsep}{3pt}
\begin{tabularx}{\linewidth}{@{}>{\raggedright\arraybackslash}X >{\centering\arraybackslash}p{.31\linewidth} >{\centering\arraybackslash}p{.20\linewidth}@{}}\toprule
Complete outcome & Probability & $\cTwoValue{X}$ (\cTwoValue{USD}) \\ \midrule
\cTwoOutcome{Cash} & $\cTwoMassHL{0.15}$ & $\cTwoValue{0}$ \\
\cTwoOutcome{Debit card} & $\cTwoMassHL{0.25}$ & $\cTwoValue{2}$ \\
\cTwoOutcome{Credit card} & $\cTwoMassHL{0.35}$ & $\cTwoValue{2}$ \\
\cTwoOutcome{Wallet} & $\cTwoMassHL{0.25}$ & $\cTwoValue{4}$ \\
\bottomrule\end{tabularx}
\end{minipage}\hfill\begin{minipage}[t]{.40\linewidth}\vspace{0pt}
\small\setlength{\tabcolsep}{3pt}
\begin{tabularx}{\linewidth}{@{}>{\centering\arraybackslash}p{.15\linewidth} >{\centering\arraybackslash}p{.19\linewidth} >{\raggedright\arraybackslash}X@{}}\toprule
$\cTwoValue{x}$ (\cTwoValue{USD}) & $\cTwoMassHL{p_X(x)}$ & Probability calculation \\ \midrule
$\cTwoValue{0}$ & $\cTwoMassHL{0.15}$ & $P(\text{Cash})=0.15$ \\
$\cTwoValue{2}$ & $\cTwoMassHL{0.60}$ & $P(\text{Debit card})$\newline ${}+P(\text{Credit card})$\newline $=0.25+0.35=0.60$ \\
$\cTwoValue{4}$ & $\cTwoMassHL{0.25}$ & $P(\text{Wallet})=0.25$ \\
\bottomrule\end{tabularx}
\end{minipage}\hfill\begin{minipage}[t]{.28\linewidth}\vspace{0pt}\centering
% Draw directly at this column's size so axis labels are not scaled down.
\begin{tikzpicture}\begin{axis}[
width=\linewidth,height=4.15cm,axis lines=left,
xmin=-1.3,xmax=5.3,ymin=0,ymax=.74,
xtick={0,2,4},minor xtick={-1,1,3,5},ytick={0,.1,.2,.3,.4,.5,.6,.7},
grid=both,grid style={gray!20},
xlabel={$\cTwoValue{x}$ (\cTwoValue{USD})},ylabel={$\cTwoMassHL{p_X(x)}$},
tick label style={font=\small},label style={font=\small},
scaled ticks=false,clip=false]
\draw[black,fill=cTwoMassColor!20] (axis cs:-.55,0) rectangle (axis cs:.55,.15);
\draw[black,fill=cTwoMassColor!20] (axis cs:1.45,0) rectangle (axis cs:2.55,.60);
\draw[black,fill=cTwoMassColor!20] (axis cs:3.45,0) rectangle (axis cs:4.55,.25);
\addplot[only marks,mark=none,nodes near coords,
 every node near coord/.append style={font=\small},
 point meta=explicit symbolic] coordinates {(0,.15) [0.15] (2,.60) [0.60] (4,.25) [0.25]};
\end{axis}\end{tikzpicture}
\end{minipage}\par
\textbf{Grouping.} Debit and credit cards both give $X=2$, so add their probabilities: $p_X(2)=P(\text{Debit card})+P(\text{Credit card})=0.60$. For $x\notin S_X$, $p_X(x)=0$.
\textbf{PMF check.} All masses are positive and $0.15+0.6+0.25=1$. The graph uses the actual numerical spacing; gaps carry no mass.
\task{2}{Accumulate in increasing support order: $\cTwoCDFHL{0.15};\ \cTwoCDFHL{0.75};\ \cTwoCDFHL{1}$. Each new height is the previous height plus the next PMF mass.}
\par\noindent\begin{minipage}[t]{.37\linewidth}\vspace{0pt}\centering
\begin{tabular}{cc}\toprule
Cutoff interval & $\cTwoCDFHL{F_X(x)}$ \\ \midrule
$x<0$ & $\cTwoCDFHL{0}$ \\
$0\leq x<2$ & $\cTwoCDFHL{0.15}$ \\
$2\leq x<4$ & $\cTwoCDFHL{0.75}$ \\
$x\geq 4$ & $\cTwoCDFHL{1}$ \\
\bottomrule\end{tabular}
\end{minipage}\hfill\begin{minipage}[t]{.60\linewidth}\vspace{0pt}\centering
\begin{tikzpicture}\begin{axis}[classCDF,height=3.5cm,xmin=-2,xmax=6,xtick={0,2,4}]
\addplot[cTwoCDFColor,thick,<-] coordinates {(-2,0)(0,0)};
\addplot[cTwoCDFColor,thick] coordinates {(0,0.15)(2,0.15)};
\addplot[cTwoCDFColor,thick] coordinates {(2,0.75)(4,0.75)};
\addplot[cTwoCDFColor,thick,->] coordinates {(4,1)(6,1)};
\addplot[only marks,mark=*,cTwoCDFColor,mark size=1.8pt,mark options={fill=white}] coordinates {(0,0)(2,0.15)(4,0.75)};
\addplot[only marks,mark=*,cTwoCDFColor,mark size=1.8pt] coordinates {(0,0.15)(2,0.75)(4,1)};
\end{axis}\end{tikzpicture}
\end{minipage}\par
\textbf{CDF check.} All heights lie in $[0,1]$ and never decrease. Filled left/open right endpoints make the graph right-continuous. Horizontal tails give limits 0 and 1; the CDF reaches 1 at 4 USD. The jumps, $0.15,\ 0.6,\ 0.25$, reproduce the PMF.
\textbf{Interpretation.} $F_X(2)=0.75$: the probability that the fee is at most 2 USD (cash or either card) is 0.75.
\endgroup
\vfill

\stageheading{Guided Example}
\problemtitle{2}{Kayak returns}{B}
\textbf{Late fees.} A rental service records the return of one kayak. The exhaustive categories are on time or early, 1--5 whole minutes late, 6--15 whole minutes late, and more than 15 minutes late, with probabilities .25, .20, .30, .25, respectively. The late fee is 0 USD, 4 USD, 4 USD and 8 USD, respectively. No other charges apply. Record the late fee.
\task{1}{Define $X$, its units and support; write the PMF table and represent it graphically. Show the outcome mapping and probability calculations, then verify the PMF.}
\textbf{Starter.} Start with on time $\mapsto0$. Apply the fee schedule, rather than using minutes as $X$.
\worklines{4}
\par\noindent\begin{minipage}[t]{.37\linewidth}\vspace{0pt}\centering
\blanktable{$x$ (USD) & $p_X(x)$}\end{minipage}\hfill\begin{minipage}[t]{.60\linewidth}\vspace{0pt}\centering
\blankgraph{$p_X(x)$}
\end{minipage}\par
\textbf{PMF checks:} \hrulefill\par
\task{2}{Write the CDF table for every real cutoff and represent it graphically. Verify its defining properties and that its jumps reproduce the PMF. Interpret $F_X(4)$ in this context.}
\textbf{Prompt.} Start below the minimum at height 0; add one grouped mass at a time. Include both tails and distinguish open and filled endpoints.
\par\noindent\begin{minipage}[t]{.37\linewidth}\vspace{0pt}\centering
\blanktable{Cutoff interval & $F_X(x)$}\end{minipage}\hfill\begin{minipage}[t]{.60\linewidth}\vspace{0pt}\centering
\blankgraph{$F_X(x)$}
\end{minipage}\par
\textbf{CDF checks and contextual interpretation:}\worklines{3}
\vfill
\clearpage
\stageheading{Independent Practice}
\problemtitle{3}{Workshop scrap buyback}{C}
\textbf{Uneven buyback payments.} A workshop scrap centre buys one standard-sized bag containing only one material type. The exhaustive types are rejected material, paper, cardboard, plastic, steel and aluminium, with probabilities .20, .12, .18, .15, .10, .25, respectively. Payments are 0 USD for rejected material, 4 USD for either paper or cardboard, 9 USD for either plastic or steel, and 15 USD for aluminium. There are no other payments. Record the buyback payment.
\task{1}{Define $X$, its units and support; write the PMF table and represent it graphically. Show the outcome mapping and probability calculations, then verify the PMF.}
\textbf{Calculation area:} include every complete outcome and its probability before grouping.
\worklines{4}
\par\noindent\begin{minipage}[t]{.37\linewidth}\vspace{0pt}\centering
\blanktable{$x$ (USD) & $p_X(x)$}\end{minipage}\hfill\begin{minipage}[t]{.60\linewidth}\vspace{0pt}\centering
\blankgraph{$p_X(x)$}
\end{minipage}\par
\textbf{PMF checks:} \hrulefill\par
\task{2}{Write the CDF table for every real cutoff and represent it graphically. Verify its defining properties and that its jumps reproduce the PMF. Interpret $F_X(9)$ in this context.}
\par\noindent\begin{minipage}[t]{.37\linewidth}\vspace{0pt}\centering
\blanktable{Cutoff interval & $F_X(x)$}\end{minipage}\hfill\begin{minipage}[t]{.60\linewidth}\vspace{0pt}\centering
\blankgraph{$F_X(x)$}
\end{minipage}\par
\textbf{CDF checks and contextual interpretation:}\worklines{3}
\vfill
\clearpage
\problemtitle{4}{Delivery loss claims}{D}
\textbf{Missing probability and capped payout.} A delivery insurer can receive only 0, 1, 2, 3 or 4 claims in one day. Their probabilities, in that order, are .15, .20, a missing probability, .20, .10. It pays 5 USD per claim, capped at 15 USD in total for the day; it pays the smaller of 5 times the count and 15 USD. There are no other payments. Record the daily payout.
\task{1}{Define $X$, its units and support; write the PMF table and represent it graphically. Show the outcome mapping and probability calculations, then verify the PMF.}
\textbf{Calculation area:} include every complete outcome and its probability before grouping.
\worklines{4}
\par\noindent\begin{minipage}[t]{.37\linewidth}\vspace{0pt}\centering
\blanktable{$x$ (USD) & $p_X(x)$}\end{minipage}\hfill\begin{minipage}[t]{.60\linewidth}\vspace{0pt}\centering
\blankgraph{$p_X(x)$}
\end{minipage}\par
\textbf{PMF checks:} \hrulefill\par
\task{2}{Write the CDF table for every real cutoff and represent it graphically. Verify its defining properties and that its jumps reproduce the PMF. Interpret $F_X(10)$ in this context.}
\par\noindent\begin{minipage}[t]{.37\linewidth}\vspace{0pt}\centering
\blanktable{Cutoff interval & $F_X(x)$}\end{minipage}\hfill\begin{minipage}[t]{.60\linewidth}\vspace{0pt}\centering
\blankgraph{$F_X(x)$}
\end{minipage}\par
\textbf{CDF checks and contextual interpretation:}\worklines{3}
\vfill
\clearpage
\problemtitle{5}{Tote bag sales}{E}
\textbf{Independent purchases and discount.} A school enterprise offers one tote bag to each of three visitors. Independently, each buys one with probability 0.3 or passes with probability 0.7, regardless of the discount. Each item costs 7 USD. Only if all three buy, subtract a 5 USD discount from the combined bill. Otherwise there is no discount. There are no other sales or charges. Record total revenue.
\task{1}{Define $X$, its units and support; write the PMF table and represent it graphically. Show the outcome mapping and probability calculations, then verify the PMF.}
\textbf{Calculation area:} include every complete outcome and its probability before grouping.
\par\noindent\begin{tabularx}{\linewidth}{|Y|Y|Y|Y|Y|Y|}\hline Outcome & Probability & $X$ (USD) & Outcome & Probability & $X$ (USD)\\\hline
\rule{0pt}{.63cm}& & & & & \\\hline
\rule{0pt}{.63cm}& & & & & \\\hline
\rule{0pt}{.63cm}& & & & & \\\hline
\rule{0pt}{.63cm}& & & & & \\\hline
\end{tabularx}
\par\noindent\begin{minipage}[t]{.37\linewidth}\vspace{0pt}\centering
\blanktable{$x$ (USD) & $p_X(x)$}\end{minipage}\hfill\begin{minipage}[t]{.60\linewidth}\vspace{0pt}\centering
\blankgraph{$p_X(x)$}
\end{minipage}\par
\textbf{PMF checks:} \hrulefill\par
\task{2}{Write the CDF table for every real cutoff and represent it graphically. Verify its defining properties and that its jumps reproduce the PMF. Interpret $F_X(14)$ in this context.}
\par\noindent\begin{minipage}[t]{.37\linewidth}\vspace{0pt}\centering
\blanktable{Cutoff interval & $F_X(x)$}\end{minipage}\hfill\begin{minipage}[t]{.60\linewidth}\vspace{0pt}\centering
\blankgraph{$F_X(x)$}
\end{minipage}\par
\textbf{CDF checks and contextual interpretation:}\worklines{3}
\vfill
\clearpage
\problemtitle{6}{Two-stage media fund}{F}
\textbf{Signed two-stage net results.} A hypothetical classroom fund receives two cash returns. The first is 17 USD, 0 USD or 7 USD with probabilities .20, .30, .50, respectively. Independently, the second is 0 USD with probability 0.4 or 7 USD with probability 0.6. These are all possible returns. A fixed participation cost of 7 USD is charged once. There are no other receipts or charges. Record the sum of both returns minus the cost; a negative net result means a loss.
\task{1}{Define $X$, its units and support; write the PMF table and represent it graphically. Show the outcome mapping and probability calculations, then verify the PMF.}
\textbf{Calculation area:} include every complete outcome and its probability before grouping.
\par\noindent\begin{tabularx}{\linewidth}{|Y|Y|Y|Y|Y|Y|}\hline Outcome & Probability & $X$ (USD) & Outcome & Probability & $X$ (USD)\\\hline
\rule{0pt}{.63cm}& & & & & \\\hline
\rule{0pt}{.63cm}& & & & & \\\hline
\rule{0pt}{.63cm}& & & & & \\\hline
\end{tabularx}
\par\noindent\begin{minipage}[t]{.37\linewidth}\vspace{0pt}\centering
\blanktable{$x$ (USD) & $p_X(x)$}\end{minipage}\hfill\begin{minipage}[t]{.60\linewidth}\vspace{0pt}\centering
\blankgraph{$p_X(x)$}
\end{minipage}\par
\textbf{PMF checks:} \hrulefill\par
\task{2}{Write the CDF table for every real cutoff and represent it graphically. Verify its defining properties and that its jumps reproduce the PMF. Interpret $F_X(0)$ in this context.}
\par\noindent\begin{minipage}[t]{.37\linewidth}\vspace{0pt}\centering
\blanktable{Cutoff interval & $F_X(x)$}\end{minipage}\hfill\begin{minipage}[t]{.60\linewidth}\vspace{0pt}\centering
\blankgraph{$F_X(x)$}
\end{minipage}\par
\textbf{CDF checks and contextual interpretation:}\worklines{3}
\vfill
\clearpage
\ifprintsolutions
\clearpage
\stageheading{Solutions}
\problemtitle{2}{Kayak returns}{B}
\textbf{Late fees.} A rental service records the return of one kayak. The exhaustive categories are on time or early, 1--5 whole minutes late, 6--15 whole minutes late, and more than 15 minutes late, with probabilities .25, .20, .30, .25, respectively. The late fee is 0 USD, 4 USD, 4 USD and 8 USD, respectively. No other charges apply. Record the late fee.
\task{1}{Let $\cTwoValue{X}$ be the late fee for one return, in USD. The two shorter delays both map to 4 USD. Their mutually exclusive masses add: $p_X(4)=.20+.30=0.5$. On-time returns map to 0 and the longest delay maps to 8.}
\begin{center}
\begin{tabular}{ccc}\toprule
Complete outcome & Probability & $\cTwoValue{X}$ (\cTwoValue{USD}) \\ \midrule
\cTwoOutcome{On time/early} & $\cTwoMassHL{0.25}$ & $\cTwoValue{0}$ \\
\cTwoOutcome{1--5 min} & $\cTwoMassHL{0.2}$ & $\cTwoValue{4}$ \\
\cTwoOutcome{6--15 min} & $\cTwoMassHL{0.3}$ & $\cTwoValue{4}$ \\
\cTwoOutcome{Over 15 min} & $\cTwoMassHL{0.25}$ & $\cTwoValue{8}$ \\
\bottomrule\end{tabular}
\end{center}
$\cTwoValue{S_X}=\cTwoValue{\{0,4,8\}}$; $p_X(x)=0$ elsewhere. The grouped PMF is:
\par\noindent\begin{minipage}[t]{.37\linewidth}\vspace{0pt}\centering
\begin{tabular}{cc}\toprule
$\cTwoValue{x}$ (\cTwoValue{USD}) & $\cTwoMassHL{p_X(x)}$ \\ \midrule
$\cTwoValue{0}$ & $\cTwoMassHL{0.25}$ \\
$\cTwoValue{4}$ & $\cTwoMassHL{0.5}$ \\
$\cTwoValue{8}$ & $\cTwoMassHL{0.25}$ \\
\bottomrule\end{tabular}
\end{minipage}\hfill\begin{minipage}[t]{.60\linewidth}\vspace{0pt}\centering
\GTenPMFGraph[color=cTwoMassColor!20,xhighlight=GTenHLValue,yhighlight=cTwoMassMarker,width=\linewidth,height=4.15cm,xlabel={$\cTwoValue{x}$ (\cTwoValue{USD})},ylabel={$\cTwoMassHL{p_X(x)}$},xstep=1]{0,4,8}{0.25,0.5,0.25}
\end{minipage}\par
\textbf{PMF check.} All masses are positive and $0.25+0.5+0.25=1$. The graph uses the actual numerical spacing; gaps carry no mass.
\task{2}{Accumulate in increasing support order: $0.25;\ 0.75;\ 1$. Each new height is the previous height plus the next PMF mass.}
\par\noindent\begin{minipage}[t]{.37\linewidth}\vspace{0pt}\centering
\begin{tabular}{cc}\toprule
Cutoff interval & $\cTwoCDFHL{F_X(x)}$ \\ \midrule
$x<0$ & $\cTwoCDFHL{0}$ \\
$0\leq x<4$ & $\cTwoCDFHL{0.25}$ \\
$4\leq x<8$ & $\cTwoCDFHL{0.75}$ \\
$x\geq 8$ & $\cTwoCDFHL{1}$ \\
\bottomrule\end{tabular}
\end{minipage}\hfill\begin{minipage}[t]{.60\linewidth}\vspace{0pt}\centering
\begin{tikzpicture}\begin{axis}[classCDF,xmin=-4,xmax=12,xtick={0,4,8}]
\addplot[cTwoCDFColor,thick,<-] coordinates {(-4,0)(0,0)};
\addplot[cTwoCDFColor,thick] coordinates {(0,0.25)(4,0.25)};
\addplot[cTwoCDFColor,thick] coordinates {(4,0.75)(8,0.75)};
\addplot[cTwoCDFColor,thick,->] coordinates {(8,1)(12,1)};
\addplot[only marks,mark=*,cTwoCDFColor,mark size=1.8pt,mark options={fill=white}] coordinates {(0,0)(4,0.25)(8,0.75)};
\addplot[only marks,mark=*,cTwoCDFColor,mark size=1.8pt] coordinates {(0,0.25)(4,0.75)(8,1)};
\end{axis}\end{tikzpicture}
\end{minipage}\par
\textbf{CDF check.} All heights lie in $[0,1]$ and never decrease. Filled left/open right endpoints make the graph right-continuous. Horizontal tails give limits 0 and 1; the CDF reaches 1 at 8 USD. The jumps, $0.25,\ 0.5,\ 0.25$, reproduce the PMF.
\textbf{Interpretation.} $F_X(4)=0.75$: the probability that the late fee is at most 4 USD (the return is no more than 15 minutes late) is 0.75.
\vfill
\clearpage
\problemtitle{3}{Workshop scrap buyback}{C}
\textbf{Uneven buyback payments.} A workshop scrap centre buys one standard-sized bag containing only one material type. The exhaustive types are rejected material, paper, cardboard, plastic, steel and aluminium, with probabilities .20, .12, .18, .15, .10, .25, respectively. Payments are 0 USD for rejected material, 4 USD for either paper or cardboard, 9 USD for either plastic or steel, and 15 USD for aluminium. There are no other payments. Record the buyback payment.
\task{1}{Let $\cTwoValue{X}$ be the buyback payment for one bag, in USD. Paper and cardboard share 4 USD, so their probabilities add to 0.3. Plastic and steel share 9 USD, giving 0.25. Keep the actual gaps 4, 5 and 6 USD on both graphs.}
\begin{center}
\begin{tabular}{ccc}\toprule
Complete outcome & Probability & $\cTwoValue{X}$ (\cTwoValue{USD}) \\ \midrule
\cTwoOutcome{Rejected} & $\cTwoMassHL{0.2}$ & $\cTwoValue{0}$ \\
\cTwoOutcome{Paper} & $\cTwoMassHL{0.12}$ & $\cTwoValue{4}$ \\
\cTwoOutcome{Cardboard} & $\cTwoMassHL{0.18}$ & $\cTwoValue{4}$ \\
\cTwoOutcome{Plastic} & $\cTwoMassHL{0.15}$ & $\cTwoValue{9}$ \\
\cTwoOutcome{Steel} & $\cTwoMassHL{0.1}$ & $\cTwoValue{9}$ \\
\cTwoOutcome{Aluminium} & $\cTwoMassHL{0.25}$ & $\cTwoValue{15}$ \\
\bottomrule\end{tabular}
\end{center}
$\cTwoValue{S_X}=\cTwoValue{\{0,4,9,15\}}$; $p_X(x)=0$ elsewhere. The grouped PMF is:
\par\noindent\begin{minipage}[t]{.37\linewidth}\vspace{0pt}\centering
\begin{tabular}{cc}\toprule
$\cTwoValue{x}$ (\cTwoValue{USD}) & $\cTwoMassHL{p_X(x)}$ \\ \midrule
$\cTwoValue{0}$ & $\cTwoMassHL{0.2}$ \\
$\cTwoValue{4}$ & $\cTwoMassHL{0.3}$ \\
$\cTwoValue{9}$ & $\cTwoMassHL{0.25}$ \\
$\cTwoValue{15}$ & $\cTwoMassHL{0.25}$ \\
\bottomrule\end{tabular}
\end{minipage}\hfill\begin{minipage}[t]{.60\linewidth}\vspace{0pt}\centering
\GTenPMFGraph[color=cTwoMassColor!20,xhighlight=GTenHLValue,yhighlight=cTwoMassMarker,width=\linewidth,height=4.15cm,xlabel={$\cTwoValue{x}$ (\cTwoValue{USD})},ylabel={$\cTwoMassHL{p_X(x)}$},xstep=1]{0,4,9,15}{0.2,0.3,0.25,0.25}
\end{minipage}\par
\textbf{PMF check.} All masses are positive and $0.2+0.3+0.25+0.25=1$. The graph uses the actual numerical spacing; gaps carry no mass.
\task{2}{Accumulate in increasing support order: $0.2;\ 0.5;\ 0.75;\ 1$. Each new height is the previous height plus the next PMF mass.}
\par\noindent\begin{minipage}[t]{.37\linewidth}\vspace{0pt}\centering
\begin{tabular}{cc}\toprule
Cutoff interval & $\cTwoCDFHL{F_X(x)}$ \\ \midrule
$x<0$ & $\cTwoCDFHL{0}$ \\
$0\leq x<4$ & $\cTwoCDFHL{0.2}$ \\
$4\leq x<9$ & $\cTwoCDFHL{0.5}$ \\
$9\leq x<15$ & $\cTwoCDFHL{0.75}$ \\
$x\geq 15$ & $\cTwoCDFHL{1}$ \\
\bottomrule\end{tabular}
\end{minipage}\hfill\begin{minipage}[t]{.60\linewidth}\vspace{0pt}\centering
\begin{tikzpicture}\begin{axis}[classCDF,xmin=-4,xmax=19,xtick={0,4,9,15}]
\addplot[cTwoCDFColor,thick,<-] coordinates {(-4,0)(0,0)};
\addplot[cTwoCDFColor,thick] coordinates {(0,0.2)(4,0.2)};
\addplot[cTwoCDFColor,thick] coordinates {(4,0.5)(9,0.5)};
\addplot[cTwoCDFColor,thick] coordinates {(9,0.75)(15,0.75)};
\addplot[cTwoCDFColor,thick,->] coordinates {(15,1)(19,1)};
\addplot[only marks,mark=*,cTwoCDFColor,mark size=1.8pt,mark options={fill=white}] coordinates {(0,0)(4,0.2)(9,0.5)(15,0.75)};
\addplot[only marks,mark=*,cTwoCDFColor,mark size=1.8pt] coordinates {(0,0.2)(4,0.5)(9,0.75)(15,1)};
\end{axis}\end{tikzpicture}
\end{minipage}\par
\textbf{CDF check.} All heights lie in $[0,1]$ and never decrease. Filled left/open right endpoints make the graph right-continuous. Horizontal tails give limits 0 and 1; the CDF reaches 1 at 15 USD. The jumps, $0.2,\ 0.3,\ 0.25,\ 0.25$, reproduce the PMF.
\textbf{Interpretation.} $F_X(9)=0.75$: the probability that the buyback payment is at most 9 USD (every type except aluminium) is 0.75.
\vfill
\clearpage
\problemtitle{4}{Delivery loss claims}{D}
\textbf{Missing probability and capped payout.} A delivery insurer can receive only 0, 1, 2, 3 or 4 claims in one day. Their probabilities, in that order, are .15, .20, a missing probability, .20, .10. It pays 5 USD per claim, capped at 15 USD in total for the day; it pays the smaller of 5 times the count and 15 USD. There are no other payments. Record the daily payout.
\task{1}{Let $\cTwoValue{X}$ be the total payout for one day, in USD. The missing probability is $1-(.15+.20+.20+.10)=.35$. Apply $X=\min(5N,15)$ to the claim count $N$: payouts are $0,5,10,15,15$. Four claims would pay 20 USD without the cap. Group the last two categories.}
\begin{center}
\begin{tabular}{ccc}\toprule
Complete outcome & Probability & $\cTwoValue{X}$ (\cTwoValue{USD}) \\ \midrule
\cTwoOutcome{0 claims} & $\cTwoMassHL{0.15}$ & $\cTwoValue{0}$ \\
\cTwoOutcome{1 claims} & $\cTwoMassHL{0.2}$ & $\cTwoValue{5}$ \\
\cTwoOutcome{2 claims} & $\cTwoMassHL{0.35}$ & $\cTwoValue{10}$ \\
\cTwoOutcome{3 claims} & $\cTwoMassHL{0.2}$ & $\cTwoValue{15}$ \\
\cTwoOutcome{4 claims} & $\cTwoMassHL{0.1}$ & $\cTwoValue{15}$ \\
\bottomrule\end{tabular}
\end{center}
$\cTwoValue{S_X}=\cTwoValue{\{0,5,10,15\}}$; $p_X(x)=0$ elsewhere. The grouped PMF is:
\par\noindent\begin{minipage}[t]{.37\linewidth}\vspace{0pt}\centering
\begin{tabular}{cc}\toprule
$\cTwoValue{x}$ (\cTwoValue{USD}) & $\cTwoMassHL{p_X(x)}$ \\ \midrule
$\cTwoValue{0}$ & $\cTwoMassHL{0.15}$ \\
$\cTwoValue{5}$ & $\cTwoMassHL{0.2}$ \\
$\cTwoValue{10}$ & $\cTwoMassHL{0.35}$ \\
$\cTwoValue{15}$ & $\cTwoMassHL{0.3}$ \\
\bottomrule\end{tabular}
\end{minipage}\hfill\begin{minipage}[t]{.60\linewidth}\vspace{0pt}\centering
\GTenPMFGraph[color=cTwoMassColor!20,xhighlight=GTenHLValue,yhighlight=cTwoMassMarker,width=\linewidth,height=4.15cm,xlabel={$\cTwoValue{x}$ (\cTwoValue{USD})},ylabel={$\cTwoMassHL{p_X(x)}$},xstep=1]{0,5,10,15}{0.15,0.2,0.35,0.3}
\end{minipage}\par
\textbf{PMF check.} All masses are positive and $0.15+0.2+0.35+0.3=1$. The graph uses the actual numerical spacing; gaps carry no mass.
\task{2}{Accumulate in increasing support order: $0.15;\ 0.35;\ 0.7;\ 1$. Each new height is the previous height plus the next PMF mass.}
\par\noindent\begin{minipage}[t]{.37\linewidth}\vspace{0pt}\centering
\begin{tabular}{cc}\toprule
Cutoff interval & $\cTwoCDFHL{F_X(x)}$ \\ \midrule
$x<0$ & $\cTwoCDFHL{0}$ \\
$0\leq x<5$ & $\cTwoCDFHL{0.15}$ \\
$5\leq x<10$ & $\cTwoCDFHL{0.35}$ \\
$10\leq x<15$ & $\cTwoCDFHL{0.7}$ \\
$x\geq 15$ & $\cTwoCDFHL{1}$ \\
\bottomrule\end{tabular}
\end{minipage}\hfill\begin{minipage}[t]{.60\linewidth}\vspace{0pt}\centering
\begin{tikzpicture}\begin{axis}[classCDF,xmin=-5,xmax=20,xtick={0,5,10,15}]
\addplot[cTwoCDFColor,thick,<-] coordinates {(-5,0)(0,0)};
\addplot[cTwoCDFColor,thick] coordinates {(0,0.15)(5,0.15)};
\addplot[cTwoCDFColor,thick] coordinates {(5,0.35)(10,0.35)};
\addplot[cTwoCDFColor,thick] coordinates {(10,0.7)(15,0.7)};
\addplot[cTwoCDFColor,thick,->] coordinates {(15,1)(20,1)};
\addplot[only marks,mark=*,cTwoCDFColor,mark size=1.8pt,mark options={fill=white}] coordinates {(0,0)(5,0.15)(10,0.35)(15,0.7)};
\addplot[only marks,mark=*,cTwoCDFColor,mark size=1.8pt] coordinates {(0,0.15)(5,0.35)(10,0.7)(15,1)};
\end{axis}\end{tikzpicture}
\end{minipage}\par
\textbf{CDF check.} All heights lie in $[0,1]$ and never decrease. Filled left/open right endpoints make the graph right-continuous. Horizontal tails give limits 0 and 1; the CDF reaches 1 at 15 USD. The jumps, $0.15,\ 0.2,\ 0.35,\ 0.3$, reproduce the PMF.
\textbf{Interpretation.} $F_X(10)=0.7$: the probability that the payout is at most 10 USD (at most two claims) is 0.7.
\vfill
\clearpage
\problemtitle{5}{Tote bag sales}{E}
\textbf{Independent purchases and discount.} A school enterprise offers one tote bag to each of three visitors. Independently, each buys one with probability 0.3 or passes with probability 0.7, regardless of the discount. Each item costs 7 USD. Only if all three buy, subtract a 5 USD discount from the combined bill. Otherwise there is no discount. There are no other sales or charges. Record total revenue.
\task{1}{Let $\cTwoValue{X}$ be the total revenue from the three visitors, in USD. Use B for buying and N for passing, in visitor order. Multiply within each complete sequence: e.g. $P(\mathrm{BNN})=0.3(0.7)^2=0.147$. The three one-buyer and three two-buyer outcomes each form a group. Only BBB receives the discount: $3(7)-5=16$ USD.}
\begin{center}
\begin{tabular}{ccc}\toprule
Complete outcome & Probability & $\cTwoValue{X}$ (\cTwoValue{USD}) \\ \midrule
\cTwoOutcome{NNN} & $\cTwoMassHL{0.343}$ & $\cTwoValue{0}$ \\
\cTwoOutcome{BNN} & $\cTwoMassHL{0.147}$ & $\cTwoValue{7}$ \\
\cTwoOutcome{NBN} & $\cTwoMassHL{0.147}$ & $\cTwoValue{7}$ \\
\cTwoOutcome{NNB} & $\cTwoMassHL{0.147}$ & $\cTwoValue{7}$ \\
\cTwoOutcome{BBN} & $\cTwoMassHL{0.063}$ & $\cTwoValue{14}$ \\
\cTwoOutcome{BNB} & $\cTwoMassHL{0.063}$ & $\cTwoValue{14}$ \\
\cTwoOutcome{NBB} & $\cTwoMassHL{0.063}$ & $\cTwoValue{14}$ \\
\cTwoOutcome{BBB} & $\cTwoMassHL{0.027}$ & $\cTwoValue{16}$ \\
\bottomrule\end{tabular}
\end{center}
$\cTwoValue{S_X}=\cTwoValue{\{0,7,14,16\}}$; $p_X(x)=0$ elsewhere. The grouped PMF is:
\par\noindent\begin{minipage}[t]{.37\linewidth}\vspace{0pt}\centering
\begin{tabular}{cc}\toprule
$\cTwoValue{x}$ (\cTwoValue{USD}) & $\cTwoMassHL{p_X(x)}$ \\ \midrule
$\cTwoValue{0}$ & $\cTwoMassHL{0.343}$ \\
$\cTwoValue{7}$ & $\cTwoMassHL{0.441}$ \\
$\cTwoValue{14}$ & $\cTwoMassHL{0.189}$ \\
$\cTwoValue{16}$ & $\cTwoMassHL{0.027}$ \\
\bottomrule\end{tabular}
\end{minipage}\hfill\begin{minipage}[t]{.60\linewidth}\vspace{0pt}\centering
\GTenPMFGraph[color=cTwoMassColor!20,xhighlight=GTenHLValue,yhighlight=cTwoMassMarker,width=\linewidth,height=4.15cm,xlabel={$\cTwoValue{x}$ (\cTwoValue{USD})},ylabel={$\cTwoMassHL{p_X(x)}$},xstep=1]{0,7,14,16}{0.343,0.441,0.189,0.027}
\end{minipage}\par
\textbf{PMF check.} All masses are positive and $0.343+0.441+0.189+0.027=1$. The graph uses the actual numerical spacing; gaps carry no mass.
\task{2}{Accumulate in increasing support order: $0.343;\ 0.784;\ 0.973;\ 1$. Each new height is the previous height plus the next PMF mass.}
\par\noindent\begin{minipage}[t]{.37\linewidth}\vspace{0pt}\centering
\begin{tabular}{cc}\toprule
Cutoff interval & $\cTwoCDFHL{F_X(x)}$ \\ \midrule
$x<0$ & $\cTwoCDFHL{0}$ \\
$0\leq x<7$ & $\cTwoCDFHL{0.343}$ \\
$7\leq x<14$ & $\cTwoCDFHL{0.784}$ \\
$14\leq x<16$ & $\cTwoCDFHL{0.973}$ \\
$x\geq 16$ & $\cTwoCDFHL{1}$ \\
\bottomrule\end{tabular}
\end{minipage}\hfill\begin{minipage}[t]{.60\linewidth}\vspace{0pt}\centering
\begin{tikzpicture}\begin{axis}[classCDF,xmin=-2,xmax=18,xtick={0,7,14,16}]
\addplot[cTwoCDFColor,thick,<-] coordinates {(-2,0)(0,0)};
\addplot[cTwoCDFColor,thick] coordinates {(0,0.343)(7,0.343)};
\addplot[cTwoCDFColor,thick] coordinates {(7,0.784)(14,0.784)};
\addplot[cTwoCDFColor,thick] coordinates {(14,0.973)(16,0.973)};
\addplot[cTwoCDFColor,thick,->] coordinates {(16,1)(18,1)};
\addplot[only marks,mark=*,cTwoCDFColor,mark size=1.8pt,mark options={fill=white}] coordinates {(0,0)(7,0.343)(14,0.784)(16,0.973)};
\addplot[only marks,mark=*,cTwoCDFColor,mark size=1.8pt] coordinates {(0,0.343)(7,0.784)(14,0.973)(16,1)};
\end{axis}\end{tikzpicture}
\end{minipage}\par
\textbf{CDF check.} All heights lie in $[0,1]$ and never decrease. Filled left/open right endpoints make the graph right-continuous. Horizontal tails give limits 0 and 1; the CDF reaches 1 at 16 USD. The jumps, $0.343,\ 0.441,\ 0.189,\ 0.027$, reproduce the PMF.
\textbf{Interpretation.} $F_X(14)=0.973$: the probability that revenue is at most 14 USD (not all three visitors buy) is 0.973.
\vfill
\clearpage
\problemtitle{6}{Two-stage media fund}{F}
\textbf{Signed two-stage net results.} A hypothetical classroom fund receives two cash returns. The first is 17 USD, 0 USD or 7 USD with probabilities .20, .30, .50, respectively. Independently, the second is 0 USD with probability 0.4 or 7 USD with probability 0.6. These are all possible returns. A fixed participation cost of 7 USD is charged once. There are no other receipts or charges. Record the sum of both returns minus the cost; a negative net result means a loss.
\task{1}{Let $\cTwoValue{X}$ be the signed net result after both stages and the cost, in USD. For each ordered pair, multiply independent probabilities, then compute $X=R_1+R_2-7$. For example, $(17,0)$ has probability $.20(0.4)=0.08$ and net result $17+0-7=10$. The pairs $(0,7)$ and $(7,0)$ both break even; group their masses. Subtract the cost only once and sort signed results numerically.}
\begin{center}
\begin{tabular}{ccc}\toprule
Complete outcome & Probability & $\cTwoValue{X}$ (\cTwoValue{USD}) \\ \midrule
\cTwoOutcome{$(17,0)$} & $\cTwoMassHL{0.08}$ & $\cTwoValue{10}$ \\
\cTwoOutcome{$(17,7)$} & $\cTwoMassHL{0.12}$ & $\cTwoValue{17}$ \\
\cTwoOutcome{$(0,0)$} & $\cTwoMassHL{0.12}$ & $\cTwoValue{-7}$ \\
\cTwoOutcome{$(0,7)$} & $\cTwoMassHL{0.18}$ & $\cTwoValue{0}$ \\
\cTwoOutcome{$(7,0)$} & $\cTwoMassHL{0.2}$ & $\cTwoValue{0}$ \\
\cTwoOutcome{$(7,7)$} & $\cTwoMassHL{0.3}$ & $\cTwoValue{7}$ \\
\bottomrule\end{tabular}
\end{center}
$\cTwoValue{S_X}=\cTwoValue{\{-7,0,7,10,17\}}$; $p_X(x)=0$ elsewhere. The grouped PMF is:
\par\noindent\begin{minipage}[t]{.37\linewidth}\vspace{0pt}\centering
\begin{tabular}{cc}\toprule
$\cTwoValue{x}$ (\cTwoValue{USD}) & $\cTwoMassHL{p_X(x)}$ \\ \midrule
$\cTwoValue{-7}$ & $\cTwoMassHL{0.12}$ \\
$\cTwoValue{0}$ & $\cTwoMassHL{0.38}$ \\
$\cTwoValue{7}$ & $\cTwoMassHL{0.3}$ \\
$\cTwoValue{10}$ & $\cTwoMassHL{0.08}$ \\
$\cTwoValue{17}$ & $\cTwoMassHL{0.12}$ \\
\bottomrule\end{tabular}
\end{minipage}\hfill\begin{minipage}[t]{.60\linewidth}\vspace{0pt}\centering
\GTenPMFGraph[color=cTwoMassColor!20,xhighlight=GTenHLValue,yhighlight=cTwoMassMarker,width=\linewidth,height=4.15cm,xlabel={$\cTwoValue{x}$ (\cTwoValue{USD})},ylabel={$\cTwoMassHL{p_X(x)}$},xstep=1]{-7,0,7,10,17}{0.12,0.38,0.3,0.08,0.12}
\end{minipage}\par
\textbf{PMF check.} All masses are positive and $0.12+0.38+0.3+0.08+0.12=1$. The graph uses the actual numerical spacing; gaps carry no mass.
\task{2}{Accumulate in increasing support order: $0.12;\ 0.5;\ 0.8;\ 0.88;\ 1$. Each new height is the previous height plus the next PMF mass.}
\par\noindent\begin{minipage}[t]{.37\linewidth}\vspace{0pt}\centering
\begin{tabular}{cc}\toprule
Cutoff interval & $\cTwoCDFHL{F_X(x)}$ \\ \midrule
$x<-7$ & $\cTwoCDFHL{0}$ \\
$-7\leq x<0$ & $\cTwoCDFHL{0.12}$ \\
$0\leq x<7$ & $\cTwoCDFHL{0.5}$ \\
$7\leq x<10$ & $\cTwoCDFHL{0.8}$ \\
$10\leq x<17$ & $\cTwoCDFHL{0.88}$ \\
$x\geq 17$ & $\cTwoCDFHL{1}$ \\
\bottomrule\end{tabular}
\end{minipage}\hfill\begin{minipage}[t]{.60\linewidth}\vspace{0pt}\centering
\begin{tikzpicture}\begin{axis}[classCDF,xmin=-10,xmax=20,xtick={-7,0,7,10,17}]
\addplot[cTwoCDFColor,thick,<-] coordinates {(-10,0)(-7,0)};
\addplot[cTwoCDFColor,thick] coordinates {(-7,0.12)(0,0.12)};
\addplot[cTwoCDFColor,thick] coordinates {(0,0.5)(7,0.5)};
\addplot[cTwoCDFColor,thick] coordinates {(7,0.8)(10,0.8)};
\addplot[cTwoCDFColor,thick] coordinates {(10,0.88)(17,0.88)};
\addplot[cTwoCDFColor,thick,->] coordinates {(17,1)(20,1)};
\addplot[only marks,mark=*,cTwoCDFColor,mark size=1.8pt,mark options={fill=white}] coordinates {(-7,0)(0,0.12)(7,0.5)(10,0.8)(17,0.88)};
\addplot[only marks,mark=*,cTwoCDFColor,mark size=1.8pt] coordinates {(-7,0.12)(0,0.5)(7,0.8)(10,0.88)(17,1)};
\end{axis}\end{tikzpicture}
\end{minipage}\par
\textbf{CDF check.} All heights lie in $[0,1]$ and never decrease. Filled left/open right endpoints make the graph right-continuous. Horizontal tails give limits 0 and 1; the CDF reaches 1 at 17 USD. The jumps, $0.12,\ 0.38,\ 0.3,\ 0.08,\ 0.12$, reproduce the PMF.
\textbf{Interpretation.} $F_X(0)=0.5$: the probability that the net result is nonpositive (a loss or break-even) is 0.5.
\vfill
\clearpage
\fi
\end{document}
